% Fragmento de inserción breve. Desarrollo y procedencia en la misma carpeta.
\subsubsection{Memoria aritmética del canal multiplicativo}
\label{prd-arit:seccion}

La conservación de la firma regional se prolonga a la evaluación
aritmética completa de APP. Para una posición orientada
\(x=(i,j,d)\in\{1,\ldots,9\}^2\times\{N,E,S,O\}\), sea
\(\chi(x)=(a,b,\epsilon)\), donde \(a\) es la coordenada fija,
\(b\) la coordenada móvil y \(\epsilon\) su orientación:
\[
 \chi(i,j,N)=(j,i,-1),\quad \chi(i,j,S)=(j,i,+1),\qquad
 \chi(i,j,E)=(i,j,+1),\quad \chi(i,j,O)=(i,j,-1).
\]
La etiqueta de eje \(\iota\in\{i,j\}\) conserva cuál de las
dos coordenadas originales se mueve. Escribiendo
\(\langle b\rangle_9=1+((b-1)\bmod9)\), los operadores y la
evaluación son
\begin{align}
 \overline P(a,b,\epsilon)&=(a,\langle b+\epsilon\rangle_9,\epsilon),
 &\overline H(a,b,\epsilon)&=(a,b,-\epsilon),\label{prd-arit:PH}\\
 n(a,b)&=ab,&r(a,b)&=\langle ab\rangle_9,\qquad
 q(a,b)=\frac{ab-r(a,b)}9.\label{prd-arit:rq}
\end{align}

\begin{proposition}[Representación exacta del transporte aritmético]
\label{prd-arit:minimal}
Se cumplen \(\chi P=\overline P\chi\) y
\(\chi H=\overline H\chi\). La representación \(\chi\)
tiene \(162\) estados y es la menos fina de las representaciones
estables bajo \(P,H\) que conservan el producto entero o,
equivalentemente, el par \((r,q)\). Conservar además \(\iota\)
recupera las \(324\) posiciones orientadas originales.
\end{proposition}
\begin{proof}
El avance modifica sólo \(b\), y la media vuelta sólo
\(\epsilon\), lo que prueba las dos identidades de transporte.
Las nueve observaciones \(n_k=n(\overline P^k\chi(x))\),
\(0\leq k<9\), recorren \(a,2a,\ldots,9a\). Así se recuperan
\(a=\min_k n_k\), \(b=n_0/a\) y el signo mediante
\((n_1/a-n_0/a)\bmod9\in\{1,8\}\). Toda equivalencia estable
que conserve \(n\) conserva estas nueve observaciones y, por
tanto, distingue todos los estados \((a,b,\epsilon)\).
La identidad \(n=r+9q\) prueba la equivalencia de las dos
evaluaciones. Cada fibra de \(\chi\) contiene las dos rutas
transpuestas: \((i,j,d)\mapsto(j,i,d^\top)\), con
\(N^\top=O,O^\top=N,E^\top=S,S^\top=E\). La etiqueta de eje distingue
esas dos preimágenes y permite reconstruir posición y dirección.
\end{proof}

El selector TRIT activa el producto en \(t\equiv2,5,8\pmod9\),
antes de cada avance. La media vuelta actúa al finalizar cada
\(27\) instantes. Si \(y_{t-1}\) denota el estado aritmético
previo, el incremento y su memoria acumulada quedan determinados por
\begin{equation}
 \eta_q(t)=\mathbf1_{\{t\equiv2\ (3)\}}q(y_{t-1}),\qquad
 \lambda_q(t)=\lambda_q(t-1)+\eta_q(t),\quad \lambda_q(0)=0.
 \label{prd-arit:incremento}
\end{equation}
Se cumple \(0\leq q(a,b)\leq a-1\). Para invertir una
actualización se deshace primero la media vuelta que corresponda,
después el avance, y finalmente se resta el cociente recuperado del
estado previo. El registro retiene fase, eje, orden de eventos y
memoria; \(\lambda_q\) es un cociclo aditivo de este transporte.

Cada bloque de \(27\) contiene nueve avances y recorre una vez
la coordenada móvil. Por ello, el retorno de posición y orientación
en \(54\), y de fase en \(108\), coexiste con
\begin{equation}
 \lambda_q(108N)=NQ_\times^+(a),\qquad
 Q_\times^+(a)=\frac{180a-4\sum_{b=1}^9r(a,b)}9,
 \qquad N\geq0.
 \label{prd-arit:retorno}
\end{equation}
En efecto, cada bloque suma los productos \(a,2a,\ldots,9a\);
sumar \(n=r+9q\) en los cuatro bloques da la fórmula. Si
\(a\) es primo con \(9\), los residuos suman \(45\), y
\(Q_\times^+(a)=20(a-1)\). En la ventana \(j\), contando desde
\(j=0\), el signo relativo del calendario es
\(\sigma_j=(-1)^{\lfloor j/3\rfloor}\), y la orientación
efectiva es \(\epsilon_j=\epsilon_0\sigma_j\).
La carga firmada por \(\sigma_j\) se anula; la carga ponderada
por la orientación efectiva se anula asimismo al multiplicarse por
\(\epsilon_0\). El libro conserva la orientación inicial, la carga
acumulada y sus incrementos. Estas cantidades pertenecen
al cursor multiplicativo; los lectores de una trayectoria que
alterna ambas hojas actúan sobre el registro de esa trayectoria.
